\scanpage{14}\mspage{6}
Je dis que ce diagramme est anticommutatif, i.e. les deux applications composées $\deleted{\unsure{a}}a\longrightarrow(\det M)^*$ diffèrent par le signe, ou encore sont déduites l'une de l'autre par la transposition dans $a$. En effet, a l'élément $\dot x$, la première associe $x\wedge y$, la deuxième $x\wedge(-y)$, ce qui prouve notre assertion. Donc on a les conditions équivalentes
\[
\begin{aligned}
(a\in SA_M)&\iff f\in SF_M\iff f'\in SF_M\\
&\iff x\wedge y\in\{\pm e_M\}.
\end{aligned}\tag*{(34)}
\]
(où $a=\{\dot x,\dot y\}$, $(x,y)\in\mathrm{Rep}(M)$)

Cela montre en particulier que
\[
SP_M=(S_M,SA_M,SF_M)\tag*{(35)}
\]
est encore un $2$-polyèdre triangulé, chaque arête étant contenue dans exactement deux faces. Il reste cependant à examiner la question de la connexité du contour combinatoire \deleted{\unsure{formé|fermé|formés}, d} associé
% archive sheet number: 13
